\subsection{A uniform alternate-field unit criterion}
\label{sec:minus7-general-source}

\begin{proposition}[K17-MINUS7-GENERAL-SOURCE]
\label{prop:minus7-general-source}
Let $e\geq2$ be even and let $\ell>\max(2e+1,7)$ be prime.
Put $E=\mathbf Q(\sqrt{-7})$, and let $\mathcal A_{e-1}$ be the
entire space of algebraic modular forms for $U(I_3/E)$, at the full
integral lattice level, of highest weight $(e-1,0,1-e)$.
Let $T_2$ denote the positive minuscule operator on this space and
the usual Hecke operator on $S_{2e+1}(\Gamma_1(7))$, respectively.
Suppose these full spaces have integral Hecke lattices at $\ell$
and that
\[
 \begin{split}
 \det\bigl((1+2^e+2^{2e})I-2^{e-1}T_2
                         \mid\mathcal A_{e-1}\bigr)&\in O^\times,\\
 \det\bigl((1+2^{2e})^3I-T_2^3
                         \mid S_{2e+1}(\Gamma_1(7))\bigr)&\in O^\times
 \end{split}
 \tag{GS7.1}
\]
at every coefficient place above $\ell$. Then
\[
             \mathcal S_\ell(e)=0,
             \qquad K_{2e}(\mathbf Z)\{\ell\}=0.
 \tag{GS7.2}
\]
Both determinants must concern the full spaces; the statement does
not allow selected eigensystems in their place.
Here $O$ denotes the valuation ring of a finite coefficient
extension at the specified place; enlarging that extension does
not change either unit condition.
\end{proposition}

\begin{proof}
We verify the source datum, both representation families, and the
strict Selmer conclusion for these general parameters. The finite
place arguments use exactly the cited local and global theorems
with their locators in Section~\ref{sec:minus7-source-coverage}.

Take $N=2$ and $\psi=1$. The field has discriminant $-7$ and class
number one by the Minkowski argument in Lemma~\ref{lem:minus7-source}.
The modulus is even and prime to the discriminant. The coefficient
prime is odd and divides neither the discriminant nor the class
number. The trivial character is quadratic and prime to $\ell$,
and $\psi(-1)=(-1)^e$ because $e$ is even. The inequality
$\ell-1>2e$ gives the required nondivisibility condition.
Finally $E\not\subset\mathbf Q(\mu_\ell)$, because $E$ ramifies
at $7$ whereas the cyclotomic field is unramified there. This is the
remaining source condition \texttt{eta}. Hence the actual
$\mathrm{GD}(\Omega,E,\ell,e,1)$ holds. The prime $2$ splits in
$E$, and the direct \texttt{OneCharU\_s} criterion permits geometric
Frobenius $s$ there, exactly as in Lemma~\ref{lem:minus7-source}.

For an actual $\rho\in\mathrm{Adm3U\_s}(\iota,e,2)$, retain its
\texttt{Adm3Alt}, \texttt{Th.Unl} and literal level witnesses.
Every \texttt{KillsMu 2} character is trivial. The cyclotomic twist
$r=\rho\chi_\ell^{-1}$ therefore has
\[
              \mathrm{Level}(\ell,1,(e-1,-1,-e-1),r).
\]
Continuity gives its \texttt{WitA} witness. The proved
\texttt{unlAll\_of} and the same common twist retain the required
above-$\ell$ \texttt{Unl1} hypotheses. Applying
\texttt{Keep.levelKA\_alt} now gives a regular algebraic
conjugate-self-dual cuspidal $\mathrm{GL}_3(\mathbf A_E)$ system
unramified at every finite place, including above $\ell$, with
infinitesimal character $(e,0,-e)$. Its compact highest weight is
$(e-1,0,1-e)$.

The three infinity characters are distinct and conjugate-orthogonal.
The associated simple generic parameter has trivial reduced
component group, so the cited KMSW generic pure-inner multiplicity
formula permits the compact infinity member and all relevant
finite local choices. Shin's odd-rank theorem supplies finite-place
temperedness. The central character is trivial: it is unramified
at all finite places, trivial at infinity and factors through the
trivial ideal class group of $E$.
At split places take the spherical member; at unramified inert
places take the hyperspecial member. At $7$ the integral hyperbolic
identification, the full lattice stabilizer, and the principal-series
packet construction in Lemma~\ref{lem:minus7-u3-coverage} are
independent of $e$ and $\ell$. In particular their only parameter
is the unitary unramified root $u$, with roots
$\{u,1,u^{-1}\}$; the proof treats every such $u$, including its
quadratic reducibility case. The compact-open fixed-vector argument
therefore applies unchanged to this parameter. It uses no
supercuspidal restriction on residue cardinality or depth-preservation
conjecture. These choices give a member of the entire space
$\mathcal A_{e-1}$.
The source twist and the rank-three algebraic Satake normalization
cancel their factors of $2$. Thus, if $\alpha$ is the eigenvalue of
$2^{e-1}T_2$, one has
\[
                  \iota\operatorname{tr}\rho(s)=\alpha/2^e.
 \tag{GS7.3}
\]
The full-space conjugation identification of the two minuscule
directions works in every bidegree, so either certified direction
has the same unit conclusion.

For an actual $\rho\in\mathrm{Adm2A}(\iota,e,14)$, the source
gives $\rho\simeq1\oplus W|_{G_E}$, the literal
$\mathrm{Level}(\ell,14,(e,-e),W)$, and a cuspidal weight-$2e+1$
newform attached to $R=W\chi_\ell^{-e}$. Attachment at almost all
primes identifies $R$ with its Deligne representation by the same
Chebotarev and Brauer--Nesbitt argument as in
Lemma~\ref{lem:minus7-gl2-conductor}.
Away from $\ell$, inertia is a sum of global mod-$7$ characters,
has zero monodromy, and is trivial outside $7$.
At $\ell$, the ordinary flag has unramified unit characters times
the cyclotomic powers $1$ and $\chi_\ell^{-2e}$. Its two
Weil--Deligne Frobenius slopes therefore differ by $2e\geq4$.
Full classical Saito compatibility and the exactness of the
potentially semistable functor give $N=0$, since
$N\varphi=\ell\varphi N$ would otherwise require slope difference
one. Finite inertia has trivial graded characters and is trivial
by characteristic-zero semisimplicity. Thus the conductor at
$\ell$ is zero. At the other primes, the cited full Carayol
compatibility preserves the conductor.

Write the two inertia characters at $7$ as restrictions of global
mod-$7$ characters $\kappa_1,\kappa_2$. Twist the automorphic
representation and its Galois representation by $\kappa_1^{-1}$.
This preserves cuspidality and weight and makes inertia
$1\oplus\kappa_2\kappa_1^{-1}$. Weight $2e+1$ is odd, so the
nebentypus is odd; a square twist preserves this oddness. The
ratio character is therefore nontrivial and odd, so its tame
conductor at $7$ is exactly one. The twisted primitive form lies
in the entire $S_{2e+1}(\Gamma_1(7))$. Untwisting at $2$ has value
$\zeta\in\mu_3$, since $2$ has order three modulo $7$. Hence
\[
        \iota\operatorname{tr}\rho(s)
                     =1+\zeta a_2(f)/2^e.
 \tag{GS7.4}
\]
This includes all original conductor-$49$ possibilities.

The source comparison trace is
$t=2^{-e}+1+2^e$. Multiplying the differences in (GS7.3) and
(GS7.4) by the unit $2^e$ gives, respectively,
\[
       \alpha-(1+2^e+2^{2e}),\qquad
       \zeta a_2(f)-(1+2^{2e}).
\]
Integral Hecke lattices make all eigenvalues integral. The first
determinant in (GS7.1) therefore makes every rank-three difference
a unit. The identity
\[
 \det(t_0^3I-T_2^3)
    =\prod_{\zeta^3=1}\det(t_0I-\zeta T_2),
       \qquad t_0=1+2^{2e},
\]
and the second determinant make every rank-two difference a unit
at every coefficient place. Every actual source representation
is covered by these two families. Lemma~\ref{lem:step17-unit}
therefore gives source Selmer length zero. For $\psi=1$ this is
$\mathcal S_\ell(e)$ at the bottom coefficient ring, by the
precise identification (K5); faithful coefficient extension and
the automatic unramified local condition at $2$ do not change
vanishing. Finally (K5), with $e$ even, gives (GS7.2).
\end{proof}
